\def\documentclass[sigconf]{acmart}

% Articles submitted to the Technical Papers program at the SIGGRAPH and SIGGRAPH Asia conferences must use the following
% \documentclass command:
%     \documentclass[acmtog,review,anonymous]{acmart}
% and include the \acmSubmissionID{} command with their article's submission ID as parameter.
%
%     \acmSubmissionID{123}

% rights management commands, sent to the author by ACM after rights form completion.

\copyrightyear{2025}
\acmYear{2025}
\setcopyright{cc}
\setcctype{by}
\acmConference[Craig-Boris '25]{Craig-Boris '25: Craig's test event for Boris LaTeX Code 2}{January 1--5, 2025}{TBa, TAS, USA}
\acmBooktitle{Craig-Boris '25: Craig's test event for Boris LaTeX Code 2 (Craig-Boris '25), January 1--5, 2025, TBa, TAS, USA}
\acmDOI{10.1145/3563045.3563055}
\acmISBN{/25/01}

% packages to add, making sure they are in https://authors.acm.org/proceedings/production-information/accepted-latex-packages.

\usepackage{graphicx}
\usepackage{subfig}

% Submission ID.
% Use this when submitting an article to a sponsored event. You'll receive a unique submission ID from the organizers of the event, 
% and this ID should be used as the parameter to this command.
%\acmSubmissionID{123-A56-BU3}

% reference and citation style.

\citestyle{acmauthoryear}

% end of preamble.

% start of document.

\begin{document}

% an optional teaser image, which appears on the first page of the paper / article, below the title / author information and above
% the content of the paper / article.

\begin{teaserfigure}
\centering
  \includegraphics[width=0.9\linewidth]{images/pexels-ralph-chang-260364-866398.jpg}
  \caption{Panoramic Photography of City Near Body of Water. Photo by Ralph Chang from Pixabay.}
  \Description{Panoramic Photography of City Near Body of Water.}
  \label{fig:teaser}
\end{teaserfigure}

% author information.


\author{Stephen Spencer}
\affiliation{%
  \department{Paul G. Allen School for Computer Science and Engineering}
  \institution{University of Washington}
  \city{Seattle}
  \state{WA}
  \country{USA}}
\email{spencer9@uw.edu}

\author{Stefano Milani}
\authornote{Equal contribution.}
\affiliation{%
  \department{Department of Computer, Control and Management Engineering (DIAG)}
  \institution{Sapienza University of Rome}
  \city{Rome}
  \country{Italy}}
\email{milani@diag.uniroma1.it}

\author{Francesca Cuomo}
\authornotemark[1]
\affiliation{%
  \department{Department of Information Engineering, Electronics and Telecommunications}
  \institution{University of Rome Sapienza}
  \city{Rome}
  \country{Italy}}
\email{francesca.cuomo@uniroma1.it}

\author{Ioannis Chatzigiannakis}
\affiliation{%
  \department{Department of Computer, Control and Management Engineering}
  \institution{Sapienza University of Rome}
  \city{Rome}
  \country{Italy}}
\email{ichatz@diag.uniroma1.it}

% This command provides customization of the author information in page headers. 
% It can safely be commented out, and a default author list generated for this purpose.
% A lengthy author list can overlap with the title information in the page headers.
%\renewcommand{\shortauthors}{Smith, et al.}

% Title of the work. 
% If your title is lengthy, consider using the "short title" parameter:
%  \title[short title]{full title}

\title{Enabling Edge Computing over LoRaWAN: A Device-Gateway Coordination Protocol}

% CCS Concepts. Required for any ACM work over two pages in length.
% Visit http://dl.acm.org/ccs and generate this content for your work.

\begin{CCSXML}
<ccs2012>
   <concept>
       <concept_id>10010147.10010371.10010352</concept_id>
       <concept_desc>Computing methodologies~Animation</concept_desc>
       <concept_significance>500</concept_significance>
       </concept>
 </ccs2012>
\end{CCSXML}

\ccsdesc[500]{Computing methodologies~Animation}

% Keywords. Optional.

\keywords{this, that, other}

% Abstract.

\begin{abstract}
Please note that TAPS will generate the title and author information for the generated PDF and HTML versions of this document from the information provided by the contact author on the completed rights form, instead of the title and author information provided by the author in the source materials. You still need to include that information if (a) you are preparing your source material in LaTeX and (b) you wish to use the ``authornote'' and ``authornotemark'' commands, as shown in this example.

Duis cursus placerat leo ac vehicula. Donec id dolor vel nunc maximus consectetur. Morbi pharetra, odio quis mattis gravida, orci dolor fermentum arcu, id tempus ante ante ut justo. Pellentesque fringilla orci a quam lacinia, non mollis orci blandit. Aliquam erat volutpat. Duis sit amet magna commodo, faucibus risus sed, dictum tortor. Cras porttitor facilisis libero non rutrum. Nullam dignissim velit a quam finibus, sed consectetur arcu mattis.
\end{abstract}

% Processes the title / author / affiliation information and builds the first part of the generated PDF.

\maketitle

\section{Introduction}
The outcome in Brussels today signals a tough road ahead to secure commitments on climate action not only from European nations, but from other big polluters as well. New president of the European Commission Ursula von der Leyen just released the EU’s Green Deal on December 11~\cite{Hagerup1993}, a sweeping package of policy proposals aimed at drastically reducing emissions by 2050. Her announcement followed the November declaration by the European Parliament of a ''climate emergency''- a phrase that activists with Greenpeace hung from the building where the Council met in Brussels.

\section{Another Section}
GOES-17 went up to work with GOES-16, another NOAA weather satellite that was launched in 2016. The two probes, which are part of the so-called GOES-R series, are able to scan most of the Western Hemisphere from the coast of Africa all the way to New Zealand. Their observations from 22,300 miles (almost 36,000 kilometers) above Earth are key to monitor hurricanes, droughts, wildfires, lighting, and fog. The two spacecraft also provide us with stunning views of our planet.

\subsection{Sample Subsection}
Following biochar’s recognition in the IPCC 2018 report, earlier this year Redmile-Gordon launched the society’s first trials to see how the material could improve plant growth. He estimates planting 10-20kg of biochar in your garden could offset the carbon from a five-mile return commute in a car for a month. Biochar is a form of charcoal produced when organic matter – for example wood, leaves or dead plants – is heated at high temperatures with little or no oxygen in a process called pyrolysis. The normal burning or decomposition of these materials would release large amounts of methane and carbon dioxide into the atmosphere. Instead, creating biochar traps this carbon in solid form for centuries; it becomes a carbon sink that can be buried underground.

\begin{figure}%
    \centering
    \subfloat[\centering label 1]{{\includegraphics[width=3.75cm]{images/pexels-pixabay-33791.jpg} }}%
    \qquad
    \subfloat[\centering label 2]{{\includegraphics[width=3.75cm]{images/pexels-pixabay-33791.jpg} }}%
    \caption{Figs. Image by Stefan Schweihofer from Pixabay.}%
    \label{fig:example}%
\end{figure}

\subsubsection{Sample Subsubsection}
GOES-17 went up to work with GOES-16, another NOAA weather satellite that was launched in 2016. The two probes, which are part of the so-called GOES-R series, are able to scan most of the Western Hemisphere from the coast of Africa all the way to New Zealand. Their observations from 22,300 miles (almost 36,000 kilometers) above Earth are key to monitor hurricanes, droughts, wildfires, lighting, and fog. The two spacecraft also provide us with stunning views of our planet.

The Taylor series expansion for the function $e^x$ is given by
\begin{equation}
e^x = 1 + x + \frac{x^2}{2} + \frac{x^3}{6} + \cdots = \sum_{n\geq 0} \frac{x^n}{n!}
\end{equation}

\paragraph{Sample Paragraph}
GOES-17 went up to work with GOES-16, another NOAA weather satellite that was launched in 2016. The two probes, which are part of the so-called GOES-R series, are able to scan most of the Western Hemisphere from the coast of Africa all the way to New Zealand. Their observations from 22,300 miles (almost 36,000 kilometers) above Earth are key to monitor hurricanes, droughts, wildfires, lighting, and fog. The two spacecraft also provide us with stunning views of our planet.

\section{Conclusion and Future Work}
Following biochar’s recognition in the IPCC 2018 report, earlier this year Redmile-Gordon launched the society’s first trials to see how the material could improve plant growth. He estimates planting 10-20kg of biochar in your garden could offset the carbon from a five-mile return commute in a car for a month. Biochar is a form of charcoal produced when organic matter – for example wood, leaves or dead plants – is heated at high temperatures with little or no oxygen in a process called pyrolysis. The normal burning or decomposition of these materials would release large amounts of methane and carbon dioxide into the atmosphere. Instead, creating biochar traps this carbon in solid form for centuries; it becomes a carbon sink that can be buried underground.

GOES-17 went up to work with GOES-16, another NOAA weather satellite that was launched in 2016. The two probes, which are part of the so-called GOES-R series, are able to scan most of the Western Hemisphere from the coast of Africa all the way to New Zealand. Their observations from 22,300 miles (almost 36,000 kilometers) above Earth are key to monitor hurricanes, droughts, wildfires, lighting, and fog. The two spacecraft also provide us with stunning views of our planet.

% The acknowledgments section is placed after the body of the work and before the bibliography.

\begin{acks}
Thanks to Corporate Lorem - \url{https://corporatelorem.kovah.de/} and ``lipsum.com'' - \url{https://www.lipsum.com/} for the ``Lorem Ipsum'' text.
\end{acks}

% The bibliography is placed after the body of the work, after the acknowledgments, and before the (optional) appendices.

\bibliographystyle{ACM-Reference-Format}
\bibliography{sample-base}

% If your work has an appendix, it goes here.

% If your work has a "figures-only" section, it goes here.

% The last command.

\end{document}
